\section{Introduction}
\label{sec:intro}

Classifiers deployed on chest radiographs are typically audited with attribution maps, which indicate where evidence is concentrated but not what change would alter the prediction. Visual counterfactual explanations close this gap: given a query image $x$ with prediction $f(x)$ and a target class $c$, a generator produces $x_c$ that (i) is scored at $c$ by the classifier, (ii) remains a realistic radiograph, and (iii) differs from $x$ only in a minimal, class-relevant region, preserving anatomy, foreign objects and patient identity.

Two observations motivate this work. The first is a modelling one. A fixed, hand-tuned guidance schedule may fit a heterogeneous test population poorly: if some images need only a gentle nudge across the decision boundary and others aggressive multi-pass editing, one schedule would either under-edit the latter or de-realize the former. We therefore frame per-step guidance control as a sequential decision problem and train a PPO agent to schedule it against a reward that trades classification progress against preservation and realism. The policy is conditioned on the image state and step index, so it can emit image- and step-dependent controls. The agent never touches the diffusion weights; it controls only the sampling procedure around a frozen generator, which keeps the action space small and the learned policy reusable at inference time on arbitrary query images.

The second observation is an evaluative one, and in retrospect the more important. The field's standard success criterion, Counterfactual Validity, is almost always computed with the same classifier whose gradients guided generation. It therefore cannot distinguish a genuine class-relevant edit from an adversarial perturbation tuned to one model. This is not a hypothetical concern in our setting: our own system reaches high CV while transferring poorly to classifiers that did not participate in guidance. We consider it a reporting problem for the subfield rather than a defect specific to our generator, and we address it with two companion measurements, one anchored to human annotation and one to model independence, which we recommend alongside CV.

\paragraph{Contributions.}
\begin{enumerate}
    \item PPO-CF: a PPO-scheduled sampling controller over a frozen, LoRA-adapted Stable Diffusion generator, with a soft, segmentation-gated background lock that blends the latent outside a dilated lung mask back toward the original image's latent at every denoising step (\S\ref{sec:method}).
    \item An evaluation protocol that pairs CV with a ground-truth-anchored Lesion Hit Rate and a held-out transfer ratio, together with the split-, seeding- and routing-level controls a control-arm comparison requires in order to mean anything (\S\ref{sec:evaluation}).
    \item A five-arm study of trained policies (one training run per arm), yielding two negative results (one of them the largest effect size we measure, against a mechanism designed to help) and one confounded improvement: the ensemble-guided policy beats an earlier single-classifier policy that differs from it in guidance and probably also in reward weights, $a_1$ range and anchoring (unverified; \S\ref{sec:controls}), so only a policy-free sweep bears on ensembling itself (\S\ref{sec:five-arms}--\ref{sec:negative-results}).
\end{enumerate}
